\documentclass{article}
\usepackage{amsmath, amssymb, amsthm}
\usepackage{hyperref}
\usepackage{geometry}
\geometry{margin=1in}

\title{Erd\H{o}s Problem \#141}
\date{Last updated: 2025-08-31}
\author{Source: erdosproblems.com}

\begin{document}
\maketitle

\noindent\textbf{Status:} OPEN \quad \textbf{No prize}

\noindent\textbf{Tags:} additive combinatorics, primes, arithmetic progressions

\medskip
\noindent\textbf{Problem Statement:}

Let $k\geq 3$. Are there $k$ consecutive primes in arithmetic progression?

\bigskip
\noindent\textbf{Remarks:}

Green and Tao \cite{GrTa08} have proved that there must always exist some $k$ primes in arithmetic progression, but these need not be consecutive. Erd\H{o}s called this conjecture 'completely hopeless at present'.

The existence of such progressions for small $k$ has been verified for $k\leq 10$, see the Wikipedia page. It is open, even for $k=3$, whether there are infinitely many such progressions.

See also [219].

This is discussed in problem A6 of Guy's collection \cite{Gu04}.

\bigskip
\noindent\textbf{References:}

[GrTa08] Green, Ben and Tao, Terence, The primes contain arbitrarily long arithmetic progressions. Ann. of Math. (2) (2008), 481-547.
 
 [Gu04] Guy, Richard K., Unsolved problems in number theory. (2004), xviii+437.

\bigskip
\noindent\small{Source: \url{https://www.erdosproblems.com/141}}
\end{document}
